% =====================================================================
\chapter{Soluciones de los ejercicios}
\label{ap:soluciones}
% =====================================================================

Cada solución se presenta en cinco partes: datos, planteamiento, desarrollo, resultado y comprobación. En los pasos intermedios se redondea a cuatro cifras significativas salvo que se indique otra cosa; los resultados exactos se dan como fracción cuando existe.

% ---------------------------------------------------------------------
\section{Tema 1}

\begin{solucion}{eje:1-desplazada}
\datos ventana $x=(x_1,\dots,x_T)$ y objetivo $y$; lote de forma $(B,T)$.
\planteamiento en la posición $t$ el modelo ve $x_{1:t}$ y debe predecir el token que lo sigue en el corpus (\cref{def:ventana}).
\pasos
\begin{enumerate}[label=\arabic*.]
  \item El token que sigue a $x_{1:t}$ en el corpus es $x_{t+1}$ (para $t<T$) o el siguiente al final de la ventana (para $t=T$). Definiendo $y_t$ como ese token, cada posición tiene su respuesta en el mismo índice.
  \item Por tanto $y$ se obtiene recortando el corpus un token más tarde: es $x$ desplazada una posición.
  \item Cada ventana aporta $T$ predicciones (\cref{prop:ventanas}) y el lote apila $B$ ventanas: $B\cdot T$ predicciones.
  \item Con $B=T=32$: $32\cdot 32=1024$.
\end{enumerate}
\resultado \textbf{$y_t=x_{t+1}$; un lote $(B,T)$ contiene $B\cdot T$ predicciones, $1024$ con $B=T=32$.}
\comprobacion en el \cref{ej:sal-y-sol}, las 6 ventanas con $T=4$ dan $6\cdot 4=24$ predicciones.
\end{solucion}

\begin{solucion}{eje:1-hola}
\datos texto «hola\nl»: $n=5$ caracteres (\tok{h}, \tok{o}, \tok{l}, \tok{a}, \nl); $T=2$.
\planteamiento por la \cref{prop:ventanas} hay $n-T$ ventanas, con inicios $i=1,\dots,n-T$.
\pasos
\begin{enumerate}[label=\arabic*.]
  \item Número de ventanas: $5-2=3$.
  \item Inicio 1: $x=\tok{ho}$, $y=\tok{ol}$.
  \item Inicio 2: $x=\tok{ol}$, $y=\tok{la}$.
  \item Inicio 3: $x=\tok{la}$, $y=\tok{a}$\nl.
  \item Predicciones: $3\cdot 2=6$.
  \item Un inicio en el cuarto carácter daría $x=\tok{a}$\nl\ y exigiría como último objetivo el sexto carácter, que no existe.
\end{enumerate}
\resultado \textbf{3 ventanas: $(\tok{ho},\tok{ol})$, $(\tok{ol},\tok{la})$, $(\tok{la},\tok{a}\nl)$; 6 predicciones.}
\comprobacion cada ventana usa $T+1=3$ caracteres consecutivos, y en un texto de 5 hay exactamente $5-3+1=3$ bloques de 3.
\end{solucion}

\begin{solucion}{eje:1-d2l}
\datos $n=10$ tokens $u_1,\dots,u_{10}$; $T=4$; la ventana que empieza en $i$ usa $u_i,\dots,u_{i+4}$.
\planteamiento una ventana necesita $T+1=5$ tokens consecutivos; en el corte sin solape, los trozos empiezan en $k+1$, $k+1+T$, \dots
\pasos
\begin{enumerate}[label=\arabic*.]
  \item (a) Inicios válidos: $i=1,\dots,n-T=6$; hay 6 ventanas, que se denotan $w_1,\dots,w_6$ por su inicio.
  \item (b) Tras descartar $k$ tokens quedan $10-k$; un trozo usa 4 de entrada y necesita uno más, luego caben $\lfloor(10-k-1)/4\rfloor$ trozos.
  \item $k=0$: $\lfloor 9/4\rfloor=2$ trozos ($w_1$, $w_5$). $k=1$: $\lfloor 8/4\rfloor=2$ ($w_2$, $w_6$). $k=2$: $\lfloor 7/4\rfloor=1$ ($w_3$). $k=3$: $\lfloor 6/4\rfloor=1$ ($w_4$).
  \item Probabilidades: $P(w_1)=P(w_5)=\frac14\cdot\frac12=\frac18$; $P(w_2)=P(w_6)=\frac18$; $P(w_3)=P(w_4)=\frac14\cdot 1=\frac14$.
  \item No es uniforme: $w_3$ y $w_4$ salen con el doble de probabilidad que las demás. Sí lo es el procedimiento de la \cref{def:lote}: elegir el inicio uniformemente entre los 6 posibles, con probabilidad $\frac16$ cada uno.
  \item (c) Las frases tienen longitudes distintas: para formar un lote rectangular habría que rellenar las cortas (cómputo desperdiciado) o descartar o cortar las largas. Con el corpus concatenado y un delimitador entre documentos, todas las ventanas tienen longitud $T$ y no hace falta relleno~\cite[§5.1]{cs336}.
\end{enumerate}
\resultado \textbf{(a) 6 ventanas; (b) 2, 2, 1 y 1 trozos para $k=0,1,2,3$, con probabilidades $\frac18,\frac18,\frac14,\frac14,\frac18,\frac18$ para $w_1,\dots,w_6$: no es uniforme, y sí lo es elegir el inicio al azar entre los 6; (c) las frases completas obligan a rellenar, y la concatenación lo evita.}
\comprobacion las probabilidades suman $4\cdot\frac18+2\cdot\frac14=1$. En el \cref{ej:sal-y-sol}, también con $n=10$ y $T=4$, hay 6 ventanas.
\end{solucion}

\begin{solucion}{eje:1-estable}
\datos $o=(1000,1001,999)$, objetivo el primer componente; pérdidas $(0{,}5;\allowbreak\,1{,}2;\allowbreak\,0{,}1;\allowbreak\,2{,}0)$.
\planteamiento \cref{prop:traslacion} con $z^{\star}=1001$; perplejidad por la \cref{def:ppl}.
\pasos
\begin{enumerate}[label=\arabic*.]
  \item $o-z^{\star}=(-1,0,-2)$.
  \item $e^{-1}=0{,}3679$, $e^{0}=1$, $e^{-2}=0{,}1353$; suma $1{,}503$; $\ln 1{,}503=0{,}4076$.
  \item $\ell=-(-1)+0{,}4076=1{,}408$ nats.
  \item En precisión simple el máximo representable es $\approx 3{,}4\cdot 10^{38}\approx e^{88{,}7}$; $e^{1000}$ da infinito, el cociente $\infty/\infty$ es indefinido (NaN) y la pérdida no puede calcularse.
  \item (b) Media: $(0{,}5+1{,}2+0{,}1+2{,}0)/4=3{,}8/4=0{,}95$; $\PPL=e^{0{,}95}=2{,}586$.
\end{enumerate}
\resultado \textbf{(a) $\ell=1{,}408$ nats, y el cálculo directo desborda; (b) $\PPL\approx 2{,}59$.}
\comprobacion $o$ es una permutación de $(2,1,0)+999$ y el objetivo es el logit intermedio; con los logits del \cref{ej:gradiente}, $\ell=-1+\ln 11{,}107=-1+2{,}408=1{,}408$.
\end{solucion}

\begin{solucion}{eje:1-ppl}
\datos $\mathcal{L}=0{,}7$ nats por carácter.
\planteamiento $\PPL=e^{\mathcal{L}}$; bits $=$ nats $/\ln 2$.
\pasos
\begin{enumerate}[label=\arabic*.]
  \item $\PPL=e^{0{,}7}=2{,}014$.
  \item $\ln 2=0{,}6931$; $0{,}7/0{,}6931=1{,}010$ bits por carácter.
  \item En base 2: $2^{1{,}010}=2{,}014$.
\end{enumerate}
\resultado \textbf{$\PPL\approx 2{,}01$; $1{,}01$ bits por carácter.}
\comprobacion $2^{\mathcal{L}/\ln 2}=e^{(\mathcal{L}/\ln 2)\ln 2}=e^{\mathcal{L}}$, luego la perplejidad no depende de la base.
\end{solucion}

\begin{solucion}{eje:1-cabeza-cero}
\datos capa de salida $z=Wh+b$ con $W=0$, $b=0$; $K$ tokens; lote con $BT$ objetivos y frecuencias empíricas $\hat\pi$.
\planteamiento calcular los logits, aplicar la \cref{prop:lnK} y, para el gradiente, la \cref{prop:gradiente} con $\partial z/\partial b=I$.
\pasos
\begin{enumerate}[label=\arabic*.]
  \item Para cualquier representación $h$, $z=0\cdot h+0=\mathbf 0$: los logits no dependen de los demás parámetros.
  \item Por la \cref{prop:lnK}, cada posición tiene $p=(1/K,\dots,1/K)$ y pérdida $\ln K$.
  \item La pérdida del lote es la media de $BT$ valores iguales a $\ln K$, es decir, $\ln K$.
  \item Gradiente respecto de $b$ en una posición: como $z=Wh+b$, $\partial\ell/\partial b=\nabla_z\ell=p-e_y=\tfrac1K\mathbf 1-e_y$.
  \item Promediando sobre el lote: $\nabla_b\mathcal{L}=\tfrac1K\mathbf 1-\frac{1}{BT}\sum_{b,t}e_{Y_{b,t}}=\tfrac1K\mathbf 1-\hat\pi$.
  \item Tras un paso de descenso con tasa $\eta$: $b=\eta\,(\hat\pi-\tfrac1K\mathbf 1)$. Los tokens más frecuentes que la media reciben un logit positivo y los menos frecuentes uno negativo.
\end{enumerate}
\resultado \textbf{$\mathcal{L}_0=\ln K$ exactamente ($3{,}829$ nats con $K=46$), y $\nabla_b\mathcal{L}=\tfrac1K\mathbf 1-\hat\pi$: el primer paso orienta el sesgo hacia las frecuencias de los tokens, es decir, hacia un modelo de unigramas.}
\comprobacion las componentes de $\tfrac1K\mathbf 1-\hat\pi$ suman $1-1=0$, como las de todo gradiente respecto de los logits (\cref{ej:gradiente}). El mini-GPT de la práctica empieza en $4{,}046$, por encima de $\ln 46$, porque su cabeza no se inicializa a cero.
\end{solucion}

% ---------------------------------------------------------------------
\section{Tema 2}

\begin{solucion}{eje:2-sin-mascara}
\datos modelo de dos o más bloques; máscara $\Mm$ de la \cref{def:mha}.
\planteamiento separar el entrenamiento sin máscara (fuga de la respuesta) de la generación sin máscara (desajuste entre entrenamiento y uso).
\pasos
\begin{enumerate}[label=\arabic*.]
  \item Sin máscara al entrenar, la posición $t$ puede atender a $x_{t+1}=y_t$. Una cabeza que copie la posición siguiente reduce la pérdida casi a cero sin aprender a predecir (\cref{obs:fuga}).
  \item Ese modelo falla al generar: en la última posición no existe el token siguiente, y la pista de la que dependía desaparece.
  \item Con máscara al entrenar y sin ella al generar: en el primer bloque, la última posición no cambia, porque no hay posiciones posteriores a las que atender.
  \item Sí cambian las representaciones de las posiciones anteriores, que ahora atienden también a posiciones posteriores.
  \item Desde el segundo bloque, la última posición lee esas representaciones distintas, que nunca vio durante el entrenamiento, y su predicción se degrada.
\end{enumerate}
\resultado \textbf{sin máscara al entrenar, el modelo copia la respuesta (pérdida casi nula que no generaliza); sin máscara solo al generar, el primer bloque de la última posición no cambia, pero los siguientes reciben entradas distintas de las del entrenamiento.}
\comprobacion el \cref{tab:sintomas} asocia una pérdida casi nula en entrenamiento y validación con la fuga del futuro.
\end{solucion}

\begin{solucion}{eje:2-mascara}
\datos puntuaciones $S$ por filas $(0,1,3)$, $(1,1,2)$, $(2,0,1)$; valores $v_1=(1,0)$, $v_2=(0,1)$, $v_3=(1,1)$.
\planteamiento por el \cref{lem:mascara}, la fila $t$ es la softmax de $(S_{t1},\dots,S_{tt})$; la salida en $t$ es $\sum_{j\le t}a_{tj}v_j$.
\pasos
\begin{enumerate}[label=\arabic*.]
  \item Fila 1: solo $j=1$; $a_{1\cdot}=(1,0,0)$.
  \item Fila 2: $\softmax(1,1)=(0{,}5;\allowbreak\,0{,}5)$; $a_{2\cdot}=(0{,}5;\allowbreak\,0{,}5;\allowbreak\,0)$.
  \item Fila 3: $e^{2}=7{,}389$, $e^{0}=1$, $e^{1}=2{,}718$; suma $11{,}107$; $a_{3\cdot}=(0{,}6652;\allowbreak\,0{,}0900;\allowbreak\,0{,}2447)$.
  \item Salidas: $o_1=v_1=(1,0)$; $o_2=0{,}5\,(1,0)+0{,}5\,(0,1)=(0{,}5;\allowbreak\,0{,}5)$.
  \item $o_3=0{,}6652\,(1,0)+0{,}0900\,(0,1)+0{,}2447\,(1,1)=(0{,}6652+0{,}2447;\;0{,}0900+0{,}2447)=(0{,}9100;\allowbreak\,0{,}3348)$.
  \item No influyen las entradas por encima de la diagonal: $S_{12}=1$, $S_{13}=3$ y $S_{23}=2$.
\end{enumerate}
\resultado \textbf{pesos $\begin{psmallmatrix}1&0&0\\0{,}5&0{,}5&0\\0{,}6652&0{,}0900&0{,}2447\end{psmallmatrix}$; salidas $(1,0)$, $(0{,}5;\allowbreak\,0{,}5)$ y $(0{,}9100;\allowbreak\,0{,}3348)$; las puntuaciones $S_{12}$, $S_{13}$ y $S_{23}$ se ignoran.}
\comprobacion cada fila suma~1 ($0{,}6652+0{,}0900+0{,}2447=0{,}9999$ por redondeo). La mayor puntuación de toda la matriz, $S_{13}=3$, no tiene ningún efecto, como exige la causalidad.
\end{solucion}

\begin{solucion}{eje:2-param32}
\datos $K=46$, $T=32$, $d=32$, $L=2$.
\planteamiento \cref{prop:parametros}: $N=Kd+Td+L(12d^{2}+10d)+2d+dK+K$; con atadura se resta $Kd$ (\cref{obs:atados}).
\pasos
\begin{enumerate}[label=\arabic*.]
  \item Bloque: $12\cdot 32^{2}=12\cdot 1024=12\,288$; $10\cdot 32=320$; total $12\,608$.
  \item Dos bloques: $2\cdot 12\,608=25\,216$.
  \item Embeddings: $46\cdot 32=1472$ (token) y $32\cdot 32=1024$ (posición).
  \item Normalización final: $2\cdot 32=64$. Cabeza: $32\cdot 46+46=1518$.
  \item Total: $1472+1024+25\,216+64+1518=29\,294$.
  \item Bloques: $25\,216/99\,584=0{,}2532$ de los del modelo base.
  \item Con pesos atados: $29\,294-1472=27\,822$.
\end{enumerate}
\resultado \textbf{$N=29\,294$ (o $27\,822$ con pesos atados); los bloques tienen aproximadamente la cuarta parte de parámetros que con $d=64$.}
\comprobacion el cociente $0{,}2532$ es cercano a $(32/64)^{2}=0{,}25$, porque el término dominante es $12d^{2}$; coincide con la fila $d=32$ del \cref{tab:hiper}.
\end{solucion}

\begin{solucion}{eje:2-cabezas}
\datos modelo base con $d=64$ y $n_h=4$ frente a $n_h=1$; pérdidas finales $0{,}324$ y $0{,}336$.
\planteamiento la \cref{prop:parametros} no depende de $n_h$; $d_k=d/n_h$.
\pasos
\begin{enumerate}[label=\arabic*.]
  \item (a) Ambos tienen $N=107\,694$ parámetros.
  \item (b) Con $n_h=4$, $d_k=64/4=16$ y cada bloque calcula 4 matrices de pesos $T\times T$. Con $n_h=1$, $d_k=64$ y una sola matriz.
  \item (c) Con el mismo número de parámetros, la diferencia está en la estructura: cada cabeza produce en cada posición \emph{una} combinación convexa de valores. Con una cabeza, la posición $t$ solo puede promediar; con cuatro, puede atender a la vez, por ejemplo, a la etiqueta de la línea y al carácter anterior, y combinar ambos resultados mediante $W_O$.
\end{enumerate}
\resultado \textbf{(a) 107\,694 en ambos casos; (b) 4 matrices de pesos con $d_k=16$ frente a una con $d_k=64$; (c) la mejora procede de poder atender a varias posiciones a la vez, no de tener más parámetros.}
\comprobacion con los mismos parámetros, la variante de una cabeza es más rápida en la práctica (tiempo relativo $0{,}7$): el número de parámetros no mide el coste de calcular y aplicar varias matrices de atención.
\end{solucion}

\begin{solucion}{eje:2-gpt2}
\datos $K=50\,257$, $T=1024$, $d=768$, $L=12$; sesgos en $Q$, $K$, $V$; cabeza atada a $E$ y sin sesgo.
\planteamiento adaptar la \cref{prop:parametros}: el sesgo de $Q,K,V$ añade $3d$ por bloque, y la cabeza no añade nada.
\pasos
\begin{enumerate}[label=\arabic*.]
  \item Bloque: $3d^{2}+3d$ (QKV) $+\,d^{2}+d$ (salida) $+\,8d^{2}+5d$ (MLP) $+\,4d$ (normalizaciones) $=12d^{2}+13d$.
  \item $d^{2}=589\,824$; $12d^{2}=7\,077\,888$; $13d=9984$; bloque $=7\,087\,872$.
  \item Doce bloques: $12\cdot 7\,087\,872=85\,054\,464$.
  \item Embeddings: $Kd=50\,257\cdot 768=38\,597\,376$; $Td=1024\cdot 768=786\,432$.
  \item Normalización final: $2d=1536$. Total: $85\,054\,464+38\,597\,376+786\,432+1536=124\,439\,808$.
  \item (b) Memoria: $124\,439\,808\cdot 4=497\,759\,232$ bytes $\approx 498$~MB.
  \item (c) Fracciones:
  \[
  \frac{85\,054\,464}{124\,439\,808}=0{,}683\ \text{(bloques)},\qquad \frac{38\,597\,376+786\,432}{124\,439\,808}=0{,}316\ \text{(embeddings)}.
  \]
\end{enumerate}
\resultado \textbf{$N=124\,439\,808$ ($\approx 124$ millones), $\approx 498$~MB en precisión simple; el 68~\% de los parámetros está en los bloques y el 32~\% en las tablas de embeddings, frente al 92~\% en bloques del mini-GPT.}
\comprobacion el resultado coincide con el tamaño de 124 millones con que se suele citar GPT-2 \emph{small} (la tabla~2 de~\cite{gpt2} lo redondea a 117 millones). La fracción de embeddings es alta porque $K=50\,257$ es grande frente a $d=768$; con pesos sin atar habría $38\,597\,376$ parámetros más.
\end{solucion}

\begin{solucion}{eje:2-permutacion}
\datos mini-GPT con $P=0$; permutación $\sigma$; $x^{\sigma}_t=x_{\sigma(t)}$.
\planteamiento inducción sobre las capas, como en el \cref{teo:causal}, usando que una suma no depende del orden de sus términos.
\pasos
\begin{enumerate}[label=\arabic*.]
  \item (a) Base: $h^{0}_t(x^{\sigma})=E_{x_{\sigma(t)}}=h^{0}_{\sigma(t)}(x)$; las filas de $H^{0}(x^{\sigma})$ son las de $H^{0}(x)$ permutadas.
  \item Paso: supóngase $h^{\ell-1}_t(x^{\sigma})=h^{\ell-1}_{\sigma(t)}(x)$ para todo $t$. Las operaciones fila a fila (normalización, proyecciones, MLP, residuales) conservan esta relación.
  \item Sin máscara, la cabeza $r$ da en la fila $t$ la suma $\sum_{j=1}^{T}a_{tj}(\Vm_r)_j$, con $a_{tj}\propto\exp\bigl((\Qm_r)_t\cdot(\Km_r)_j/\sqrt{d_k}\bigr)$. Para $x^{\sigma}$, la consulta es la de la posición $\sigma(t)$ de $x$, y la suma recorre las claves y valores de todas las posiciones de $x$ en otro orden; su valor coincide con el de la fila $\sigma(t)$ para $x$. Por inducción, $z_t(x^{\sigma})=z_{\sigma(t)}(x)$.
  \item (b) Con máscara y $L=1$, la posición $T$ atiende a todas las posiciones $j\le T$, es decir, a todas. Si $\sigma(T)=T$, su consulta no cambia ($h^{0}_T=E_{x_T}$), y el conjunto de pares (clave, valor) es el mismo, reordenado. La suma, y con ella $z_T$, no cambia.
  \item (c) Con $L\ge 2$, la segunda capa en la posición $T$ lee las representaciones $h^{1}_j$ de la primera, y cada $h^{1}_j$ depende del prefijo $x_{1:j}$ por la máscara. Al permutar, los prefijos cambian como conjuntos (por ejemplo, $h^{1}_1$ pasa a depender de $x_{\sigma(1)}$ en lugar de $x_1$), de modo que el conjunto $\{h^{1}_j\}$ no es una reordenación del original.
  \item La máscara rompe, por tanto, la simetría de permutación: la posición $j$ atiende a exactamente $j$ tokens, y esa asimetría permite al modelo inferir el orden sin codificación posicional. CS336 propone comprobarlo empíricamente con el experimento \emph{no\_pos\_emb} (NoPE) de su §7.3~\cite{cs336}.
\end{enumerate}
\resultado \textbf{(a) sin máscara y sin posiciones, el modelo es equivariante por permutaciones para todo $L$; (b) con máscara y un bloque, $z_T$ es invariante frente a las permutaciones que fijan $T$; (c) con dos o más bloques la máscara introduce información de orden, por lo que un decoder puede entrenarse sin codificación posicional.}
\comprobacion para $T=2$ en (a): intercambiar $x_1$ y $x_2$ intercambia $z_1$ y $z_2$, y en (b) con $T=3$, intercambiar $x_1$ y $x_2$ deja invariante $z_3$ pero no $z_2$, que pasa de ver $\{x_1,x_2\}$ con consulta $x_2$ a verlo con consulta $x_1$.
\end{solucion}

% ---------------------------------------------------------------------
\section{Tema 3}

\begin{solucion}{eje:3-orden}
\planteamiento revisar qué hace cada paso de la \cref{def:paso} y qué cantidades modifica.
\pasos
\begin{enumerate}[label=\arabic*.]
  \item El recorte necesita el gradiente completo, luego va después de la retropropagación. Debe ir antes del paso del optimizador porque AdamW incorpora $g$ a sus medias móviles: un gradiente anómalo sin recortar contaminaría $m$ y $v$ durante muchos pasos (unos $1/(1-\beta_2)$ en el caso de $v$).
  \item Con la norma global, el gradiente recortado es un múltiplo positivo del original y conserva la dirección de descenso (\cref{prop:recorte}); por tensor, la dirección cambia.
  \item El calentamiento evita pasos grandes cuando los momentos de Adam se han estimado con pocos gradientes y estos son grandes e irregulares.
  \item El decaimiento coseno reduce la tasa al final, lo que disminuye el exceso de pérdida debido al ruido del gradiente estocástico y permite afinar.
\end{enumerate}
\resultado \textbf{recorte entre \texttt{backward} y \texttt{step}, con norma global para no cambiar la dirección; el calentamiento protege el arranque y el coseno afina el final.}
\comprobacion en el \cref{tab:hiper}, quitar el calentamiento empeora la pérdida final de $0{,}297$ a $0{,}306$ y quitar el coseno, de $0{,}324$ a $0{,}337$.
\end{solucion}

\begin{solucion}{eje:3-calendario}
\datos $W=200$, $S=2000$, $\eta_{\max}=10^{-3}$, $\eta_{\min}=10^{-4}$; $s=9$ y $s=1000$.
\planteamiento \cref{def:calendario}: rama lineal si $s<W$; rama coseno si no.
\pasos
\begin{enumerate}[label=\arabic*.]
  \item $s=9<200$: $\eta_9=10^{-3}\cdot 10/200=5\cdot 10^{-5}$.
  \item $s=1000\ge 200$: progreso $(1000-200)/(2000-200)=800/1800=0{,}4444$.
  \item Ángulo $\pi\cdot 0{,}4444=1{,}396$ rad; $\cos 1{,}396=0{,}1736$.
  \item $\tfrac12(\eta_{\max}-\eta_{\min})=\tfrac12\cdot 9\cdot 10^{-4}=4{,}5\cdot 10^{-4}$; $1+0{,}1736=1{,}1736$.
  \item $\eta_{1000}=10^{-4}+4{,}5\cdot 10^{-4}\cdot 1{,}1736=10^{-4}+5{,}281\cdot 10^{-4}=6{,}281\cdot 10^{-4}$.
\end{enumerate}
\resultado \textbf{$\eta_9=5\cdot 10^{-5}$ y $\eta_{1000}\approx 6{,}28\cdot 10^{-4}$.}
\comprobacion el punto medio del coseno es $s=1100$, con $\eta=5{,}5\cdot 10^{-4}$; $s=1000$ está antes y su tasa es mayor. Por la \cref{prop:calendario}~(iv), $\eta_{1200}=1{,}1\cdot 10^{-3}-6{,}281\cdot 10^{-4}=4{,}719\cdot 10^{-4}$, que coincide con el cálculo directo.
\end{solucion}

\begin{solucion}{eje:3-cs336-cal}
\datos $\alpha_{\max}=10^{-3}$, $\alpha_{\min}=10^{-4}$, $T_w=100$, $T_c=1000$; $\alpha_{\max}-\alpha_{\min}=9\cdot 10^{-4}$.
\planteamiento aplicar cada rama de la \cref{obs:cs336-cal}; en el coseno, la fracción es $(t-100)/900$.
\pasos
\begin{enumerate}[label=\arabic*.]
  \item $t=0$: $(0/100)\cdot 10^{-3}=0$. $t=50$: $(50/100)\cdot 10^{-3}=5\cdot 10^{-4}$.
  \item $t=100$: fracción $0$, $\cos 0=1$; $\alpha=10^{-4}+\tfrac12\cdot 2\cdot 9\cdot 10^{-4}=10^{-3}$.
  \item $t=325$: fracción $0{,}25$; $\cos(\pi/4)=0{,}7071$; $\tfrac12\cdot 1{,}7071\cdot 9\cdot 10^{-4}=7{,}682\cdot 10^{-4}$; $\alpha=8{,}682\cdot 10^{-4}$.
  \item $t=550$: fracción $0{,}5$; $\cos(\pi/2)=0$; $\alpha=10^{-4}+4{,}5\cdot 10^{-4}=5{,}5\cdot 10^{-4}$.
  \item $t=775$: fracción $0{,}75$; $\cos(3\pi/4)=-0{,}7071$; $\tfrac12\cdot 0{,}2929\cdot 9\cdot 10^{-4}=1{,}318\cdot 10^{-4}$; $\alpha=2{,}318\cdot 10^{-4}$.
  \item $t=1000$: $\cos\pi=-1$; $\alpha=10^{-4}$. $t=1200>T_c$: $\alpha=\alpha_{\min}=10^{-4}$.
\end{enumerate}
\resultado \textbf{$0$; $5\cdot 10^{-4}$; $10^{-3}$; $8{,}682\cdot 10^{-4}$; $5{,}5\cdot 10^{-4}$; $2{,}318\cdot 10^{-4}$; $10^{-4}$; $10^{-4}$.} CS336 empieza en $\alpha_0=0$ y llega al pico en $t=T_w$, y mantiene $\alpha_{\min}$ después de $T_c$; el calendario del módulo empieza en $\eta_{\max}/W$, llega al pico en $s=W-1$ y termina en $s=S-1$.
\comprobacion $t=325$ y $t=775$ son simétricos respecto del punto medio: $8{,}682\cdot 10^{-4}+2{,}318\cdot 10^{-4}=1{,}1\cdot 10^{-3}=\alpha_{\max}+\alpha_{\min}$, como en la \cref{prop:calendario}~(iv).
\end{solucion}

\begin{solucion}{eje:3-clip}
\datos $c=1$, $\varepsilon=10^{-6}$.
\planteamiento calcular $\norm{g}_2$; si es al menos $c$, multiplicar por $c/(\norm{g}_2+\varepsilon)$ (\cref{def:recorte}).
\pasos
\begin{enumerate}[label=\arabic*.]
  \item (a) $\norm{(3,4,0)}_2=\sqrt{9+16}=5\ge 1$. Factor $1/(5+10^{-6})=0{,}2000$; $g'=(0{,}6;\allowbreak\,0{,}8;\allowbreak\,0)$, de norma $5/(5+10^{-6})=0{,}9999998$.
  \item (b) $\norm{(0{,}3;\allowbreak\,0{,}4;\allowbreak\,0)}_2=\sqrt{0{,}09+0{,}16}=0{,}5<1$: no se modifica.
  \item (c) Norma global: $\sqrt{3^{2}+4^{2}+12^{2}}=\sqrt{169}=13$. Factor $1/13=0{,}07692$: $(0{,}2308;\allowbreak\,0{,}3077)$ y $(0{,}9231)$.
  \item Fracción de la norma del primer tensor: antes y después del recorte global, $5/13=0{,}3846$.
  \item Por tensor: $(3,4)\to(0{,}6;\allowbreak\,0{,}8)$ con norma~1 y $(12)\to(1)$. Norma total $\sqrt 2=1{,}414$; la fracción del primer tensor pasa a $1/\sqrt2=0{,}7071$.
\end{enumerate}
\resultado \textbf{(a) $(0{,}6;\allowbreak\,0{,}8;\allowbreak\,0)$; (b) sin cambios; (c) $(0{,}2308;\allowbreak\,0{,}3077)$ y $(0{,}9231)$. El recorte global mantiene la fracción $0{,}385$ del primer tensor; el recorte por tensor la eleva a $0{,}707$ y cambia la dirección.}
\comprobacion $0{,}2308^{2}+0{,}3077^{2}+0{,}9231^{2}=0{,}0533+0{,}0947+0{,}8521=1{,}000$.
\end{solucion}

\begin{solucion}{eje:3-sgd}
\datos $L(w)=\frac1{100}\sum_i w_i^{2}$; tasa $\alpha/\sqrt{t+1}$ en el paso $t$.
\planteamiento la actualización es lineal en $w$, y $L$ es cuadrática: si todos los pesos se multiplican por $f$, $L$ se multiplica por $f^{2}$.
\pasos
\begin{enumerate}[label=\arabic*.]
  \item (a) $\partial L/\partial w_i=2w_i/100=w_i/50$. Luego $w\leftarrow w-\frac{\alpha}{\sqrt{t+1}}\cdot\frac{w}{50}=\Bigl(1-\frac{\alpha}{50\sqrt{t+1}}\Bigr)w=f_t\,w$.
  \item (b) $L_3/L_0=(f_0f_1f_2)^{2}$, con $\sqrt2=1{,}414$ y $\sqrt3=1{,}732$.
  \item $\alpha=1$: $f=(0{,}98;\allowbreak\,0{,}9859;\allowbreak\,0{,}9885)$; producto $0{,}9551$; cociente $0{,}912$.
  \item $\alpha=10$: $f=(0{,}8;\allowbreak\,0{,}8586;\allowbreak\,0{,}8845)$; producto $0{,}6075$; cociente $0{,}369$.
  \item $\alpha=100$: $f=(-1;\allowbreak\,-0{,}4142;\allowbreak\,-0{,}1547)$; producto $-0{,}06408$; cociente $0{,}0041$.
  \item $\alpha=1000$: $f=(-19;\allowbreak\,-13{,}14;\allowbreak\,-10{,}55)$; producto en valor absoluto $2634$; cociente $6{,}94\cdot 10^{6}$.
\end{enumerate}
\resultado \textbf{$L_3/L_0=0{,}912$, $0{,}369$, $0{,}0041$ y $6{,}9\cdot 10^{6}$: las tres primeras tasas hacen decrecer la pérdida (la tercera, tras un primer paso que no la cambia) y $\alpha=1000$ diverge.}
\comprobacion la pérdida decrece en el paso $t$ si y solo si $\abs{f_t}<1$, es decir, $0<\alpha<100\sqrt{t+1}$. Con $\alpha=100$, $f_0=-1$: el peso cambia de signo y $L$ no varía; con $\alpha=1000$, $\abs{f_t}>1$ en los tres pasos.
\end{solucion}

\begin{solucion}{eje:3-flops}
\datos $N=107\,694$; 2000 pasos; $B=T=32$; rendimiento efectivo $3\cdot 10^{10}$ FLOP/s.
\planteamiento $C\approx 6ND$ con $D=\text{pasos}\cdot B\cdot T$.
\pasos
\begin{enumerate}[label=\arabic*.]
  \item $D=2000\cdot 1024=2\,048\,000$ tokens.
  \item $6N=646\,164$; $C=646\,164\cdot 2\,048\,000=1{,}323\cdot 10^{12}$ FLOP.
  \item Tiempo: $1{,}323\cdot 10^{12}/(3\cdot 10^{10})=44{,}1$ s.
\end{enumerate}
\resultado \textbf{$C\approx 1{,}32\cdot 10^{12}$ FLOP; unos 44~s.}
\comprobacion son $2000/800=2{,}5$ veces los pasos del \cref{ej:coste}: $2{,}5\cdot 5{,}29\cdot 10^{11}=1{,}32\cdot 10^{12}$ y $2{,}5\cdot 17{,}6=44$~s.
\end{solucion}

\begin{solucion}{eje:3-proyeccion}
\datos $\mathcal{B}_c=\{u:\norm{u}_2\le c\}$; $g$ arbitrario.
\planteamiento acotar $\norm{g-u}_2$ inferiormente con la desigualdad triangular y encontrar el punto que alcanza la cota.
\pasos
\begin{enumerate}[label=\arabic*.]
  \item Si $\norm{g}_2\le c$, entonces $g\in\mathcal{B}_c$ y su distancia a sí mismo es~0: $\Pi_c(g)=g$.
  \item Si $\norm{g}_2>c$, para todo $u\in\mathcal{B}_c$: $\norm{g-u}_2\ge\norm{g}_2-\norm{u}_2\ge\norm{g}_2-c$.
  \item El punto $u^{\star}=c\,g/\norm{g}_2$ está en $\mathcal{B}_c$ y $\norm{g-u^{\star}}_2=(1-c/\norm{g}_2)\norm{g}_2=\norm{g}_2-c$: alcanza la cota.
  \item Unicidad: la igualdad en la primera desigualdad exige que $u$ sea un múltiplo no negativo de $g$, y en la segunda que $\norm{u}_2=c$; el único punto que cumple ambas es $u^{\star}$.
  \item Con $\varepsilon=0$, la \cref{def:recorte} da $g$ si $\norm{g}_2<c$ y $c\,g/\norm{g}_2$ si $\norm{g}_2\ge c$ (en $\norm{g}_2=c$ ambas expresiones coinciden): es $\Pi_c$.
  \item Por tensor: $\norm{g-u}_2^{2}=\sum_q\norm{g^{(q)}-u^{(q)}}_2^{2}$ es una suma de términos que dependen de bloques distintos, y cada uno se minimiza por separado proyectando $g^{(q)}$ sobre su bola.
\end{enumerate}
\resultado \textbf{el recorte global con $\varepsilon=0$ es la proyección euclídea sobre la bola $\mathcal{B}_c$; el recorte por tensor es la proyección sobre el producto cartesiano $\mathcal{B}^{(1)}_c\times\dots\times\mathcal{B}^{(P)}_c$ de una bola por tensor.}
\comprobacion con $g=(3,4,12)$ y $c=1$, $u^{\star}=(3,4,12)/13$ está a distancia $13-1=12$ de $g$, mientras que el recorte por tensor $(0{,}6;\allowbreak\,0{,}8;\allowbreak\,1)$ está a distancia $\sqrt{2{,}4^{2}+3{,}2^{2}+11^{2}}=\sqrt{137}=11{,}70$, menor porque pertenece al producto de bolas, un conjunto mayor que $\mathcal{B}_1$, pero fuera de $\mathcal{B}_1$, ya que su norma es $\sqrt 2>1$.
\end{solucion}

% ---------------------------------------------------------------------
\section{Tema 4}

\begin{solucion}{eje:4-estado}
\datos AdamW guarda $m$, $v$ y el contador $t$; el calendario depende del paso $s$.
\planteamiento analizar qué ocurre al reanudar si esos datos se pierden.
\pasos
\begin{enumerate}[label=\arabic*.]
  \item Sin el estado del optimizador, $m=v=0$ y $t=0$: el primer paso usa momentos estimados con un único gradiente.
  \item Ese paso tiene la magnitud de un paso de Adam normalizado sin la tasa reducida del calentamiento, y su dirección puede ser mala: la pérdida da un salto y tarda en recuperarse.
  \item Sin el número de paso, el calendario vuelve a empezar: se repite el calentamiento y la tasa vuelve a subir hasta $\eta_{\max}$ cuando ya debería ser pequeña.
\end{enumerate}
\resultado \textbf{sin el estado del optimizador, el primer paso tras reanudar es anómalo y la curva da un salto; sin el número de paso, el calendario se reinicia.}
\comprobacion el estado ocupa dos números por parámetro, el doble que el modelo, coherente con los $16N$ bytes de la \cref{prop:memoria} ($4N$ del modelo, $4N$ de los gradientes y $8N$ del estado).
\end{solucion}

\begin{solucion}{eje:4-fijos}
\datos pérdida estimada como media de 20 lotes de validación fijos.
\planteamiento \cref{prop:lotes-comunes}.
\pasos
\begin{enumerate}[label=\arabic*.]
  \item La pérdida de un lote es una estimación ruidosa de la pérdida esperada.
  \item Con lotes distintos en cada evaluación, la diferencia entre dos puntos de la curva mezcla el cambio del modelo con el azar del muestreo, y su varianza es $\Var(\hat{\mathcal{L}}_a)+\Var(\hat{\mathcal{L}}_b)$.
  \item Con los mismos lotes, las dos estimaciones están positivamente correlacionadas (los lotes difíciles lo son para ambos estados) y el término $-2\Cov$ reduce la varianza de la diferencia.
\end{enumerate}
\resultado \textbf{con lotes fijos, la diferencia entre evaluaciones refleja sobre todo el cambio del modelo, porque el ruido común se cancela.}
\comprobacion en el caso extremo $\theta_a=\theta_b$, con lotes comunes la diferencia es exactamente~0, mientras que con lotes independientes es una variable aleatoria de varianza $2\Var(\hat{\mathcal{L}}_a)$.
\end{solucion}

\begin{solucion}{eje:4-linea}
\datos 4 elecciones uniformes de 6 opciones; 4 tipos de línea; 58 caracteres.
\planteamiento las elecciones son uniformes e independientes y el resto de caracteres está determinado: la entropía de la línea es el logaritmo del número de posibilidades.
\pasos
\begin{enumerate}[label=\arabic*.]
  \item Posibilidades: $6^{4}=1296$; $\ln 1296=4\ln 6=4\cdot 1{,}792=7{,}167$ nats.
  \item Tipo de línea: $\ln 4=1{,}386$ nats. Total: $7{,}167+1{,}386=8{,}553$ nats.
  \item Por carácter: $8{,}553/58=0{,}1475$ nats.
\end{enumerate}
\resultado \textbf{$\approx 0{,}147$ nats por carácter.}
\comprobacion es del mismo orden que la $H\approx 0{,}125$ del corpus completo (\cref{ej:entropia-corpus}) y muy inferior a los $0{,}324$ del modelo.
\end{solucion}

\begin{solucion}{eje:4-brecha}
\datos entrenamiento $0{,}264$; validación $0{,}324$; $H\approx 0{,}125$ nats por carácter.
\planteamiento $\PPL=e^{\mathcal{L}}$; por el \cref{teo:cota}, en media, $\mathcal{L}_{\mathrm{val}}=H+\KL$.
\pasos
\begin{enumerate}[label=\arabic*.]
  \item (a) $e^{0{,}264}=1{,}302$ y $e^{0{,}324}=1{,}383$; cociente $e^{0{,}324-0{,}264}=e^{0{,}06}=1{,}062$.
  \item (b) $\KL\approx 0{,}324-0{,}125=0{,}199$ nats por carácter.
  \item (c) $e^{0{,}125}=1{,}133$.
\end{enumerate}
\resultado \textbf{(a) perplejidades $1{,}30$ y $1{,}38$, con cociente $1{,}06$; (b) $\KL\approx 0{,}20$ nats por carácter; (c) un modelo perfecto dudaría en media entre $1{,}13$ opciones por carácter: casi todos los caracteres están determinados y la incertidumbre se concentra en unas pocas elecciones por línea.}
\comprobacion la brecha de $0{,}06$ nats es pequeña frente a la distancia a la cota, $0{,}199$: el margen de mejora está más en la capacidad y el entrenamiento que en la regularización.
\end{solucion}

\begin{solucion}{eje:4-memoria}
\datos $N=609\,838$ y $N=7\cdot 10^{9}$; 4 bytes por número; GPU de 80~GB.
\planteamiento \cref{prop:memoria}: $16N$ bytes.
\pasos
\begin{enumerate}[label=\arabic*.]
  \item (a) Parámetros: $4\cdot 609\,838=2\,439\,352$ bytes $=2{,}44$~MB; lo mismo los gradientes.
  \item Estado de AdamW: $8\cdot 609\,838=4\,878\,704$ bytes $=4{,}88$~MB. Total: $16\cdot 609\,838=9\,757\,408$ bytes $\approx 9{,}76$~MB.
  \item (b) $16\cdot 7\cdot 10^{9}=1{,}12\cdot 10^{11}$ bytes $=112$~GB.
  \item GPU necesarias: $112/80=1{,}4$, luego al menos 2.
\end{enumerate}
\resultado \textbf{(a) $\approx 9{,}76$~MB; (b) $112$~GB, que exigen al menos 2 GPU de $80$~GB solo para parámetros, gradientes y estado del optimizador.}
\comprobacion $9{,}76$~MB es $609\,838/107\,694=5{,}66$ veces el valor del modelo base, $1{,}72$~MB: $5{,}66\cdot 1{,}72=9{,}75$.
\end{solucion}

\begin{solucion}{eje:4-condicional}
\datos distribución conjunta $p(c,x)$; modelo $q(x\mid c)$ positivo donde $p(c,x)>0$.
\planteamiento descomponer la esperanza por contextos y aplicar el \cref{teo:cota} a cada uno.
\pasos
\begin{enumerate}[label=\arabic*.]
  \item $\E_p[-\ln q(X\mid C)]=\sum_{c}p(c)\sum_x p(x\mid c)\bigl(-\ln q(x\mid c)\bigr)=\sum_c p(c)\,H\bigl(p(\cdot\mid c),q(\cdot\mid c)\bigr)$.
  \item Por el \cref{teo:cota}, para cada $c$ con $p(c)>0$: $H\bigl(p(\cdot\mid c),q(\cdot\mid c)\bigr)=H\bigl(p(\cdot\mid c)\bigr)+\KL\bigl(p(\cdot\mid c)\Vert q(\cdot\mid c)\bigr)$.
  \item Sumando: $\E_p[-\ln q(X\mid C)]=H(X\mid C)+\sum_c p(c)\,\KL\bigl(p(\cdot\mid c)\Vert q(\cdot\mid c)\bigr)$, porque $\sum_c p(c)H(p(\cdot\mid c))=H(X\mid C)$.
  \item Cada divergencia es no negativa y el peso $p(c)$ también, luego la suma es no negativa: se obtiene la desigualdad.
  \item Igualdad si y solo si cada término es nulo, es decir, $\KL\bigl(p(\cdot\mid c)\Vert q(\cdot\mid c)\bigr)=0$ para todo $c$ con $p(c)>0$, lo que ocurre si y solo si $q(\cdot\mid c)=p(\cdot\mid c)$.
  \item En un modelo de lenguaje, $C$ es el contexto $x_{1:t}$ y $X$ el token siguiente; la pérdida esperada de la \cref{def:perdida} es exactamente $\E_p[-\ln q(X\mid C)]$.
\end{enumerate}
\resultado \textbf{$\E_p[-\ln q(X\mid C)]=H(X\mid C)+\E_{C}\bigl[\KL\bigl(p(\cdot\mid C)\Vert q(\cdot\mid C)\bigr)\bigr]\ge H(X\mid C)$, con igualdad si y solo si el modelo coincide con la distribución condicionada real en todo contexto posible.}
\comprobacion si $C$ es constante, $H(X\mid C)=H(X)$ y se recupera el \cref{teo:cota}.
\end{solucion}

% ---------------------------------------------------------------------
\section{Tema 5}

\begin{solucion}{eje:5-conceptos}
\planteamiento aplicar las \cref{def:topk,def:topp} y la \cref{prop:limites}.
\pasos
\begin{enumerate}[label=(\alph*)]
  \item $k=1$: solo sobrevive el mayor logit; se elige siempre el token más probable (decodificación voraz).
  \item $k=K$: no se elimina ningún logit; es el muestreo del modelo.
  \item $\rho=1$: $V(1)$ es el menor conjunto de masa~1, el de los tokens con probabilidad positiva, que con la softmax son todos; no cambia nada.
  \item $\tau\to 0^{+}$: la distribución tiende al vector \emph{one-hot} del mayor logit, como la voraz.
  \item $\tau=3$: la distribución se acerca a la uniforme y los 46 caracteres, incluidos dígitos y signos, salen con probabilidades parecidas; el texto se degrada.
\end{enumerate}
\resultado \textbf{(a) y (d) son deterministas (salvo empates); (b) y (c) equivalen a muestrear del modelo; (e) aproxima el muestreo uniforme.}
\comprobacion por el \cref{teo:entropia-tau}, la entropía con $\tau=3$ es mayor que con $\tau=1$, y está acotada por $\ln 46=3{,}829$.
\end{solucion}

\begin{solucion}{eje:5-topk}
\datos $z=(3,1,1,0)$; $k=2$; $\tau=1$.
\planteamiento quedarse con los índices de $I_2(z)$, poner el resto a $-\infty$ y aplicar la softmax.
\pasos
\begin{enumerate}[label=\arabic*.]
  \item Ordenados: $3,1,1,0$; el segundo y el tercero empatan.
  \item Convenio de la \cref{def:topk} (empate al índice menor): $I_2=\{1,2\}$. $e^{3}=20{,}09$, $e^{1}=2{,}718$; suma $22{,}80$; $p=(0{,}8808;\allowbreak\,0{,}1192;\allowbreak\,0;\allowbreak\,0)$.
  \item Convenio alternativo (todo logit $\ge 1$): sobreviven los índices 1, 2 y 3. Suma $20{,}09+2\cdot 2{,}718=25{,}52$; $p=(0{,}7870;\allowbreak\,0{,}1065;\allowbreak\,0{,}1065;\allowbreak\,0)$.
\end{enumerate}
\resultado \textbf{$(0{,}881;\allowbreak\,0{,}119;\allowbreak\,0;\allowbreak\,0)$ con desempate por índice y $(0{,}787;\allowbreak\,0{,}107;\allowbreak\,0{,}107;\allowbreak\,0)$ conservando los empatados.}
\comprobacion ambas suman~1, y la primera coincide con la que devuelve una implementación que elige exactamente $k$ índices. Con logits reales, los empates exactos son raros.
\end{solucion}

\begin{solucion}{eje:5-temp}
\datos $z=(3,1,0)$; $\tau=0{,}5$, $1$, $2$.
\planteamiento $p=\softmax(z/\tau)$; $H=-\sum p_i\ln p_i$; $\Var_p(z)=\sum p_iz_i^{2}-(\sum p_iz_i)^{2}$.
\pasos
\begin{enumerate}[label=\arabic*.]
  \item $\tau=0{,}5$: $z/\tau=(6,2,0)$; $e^{6}=403{,}4$, $e^{2}=7{,}389$, $1$; suma $411{,}8$; $p=(0{,}9796;\allowbreak\,0{,}0179;\allowbreak\,0{,}0024)$; $H=0{,}0202+0{,}0721+0{,}0146=0{,}1069$.
  \item $\tau=1$: $e^{3}=20{,}09$, $e^{1}=2{,}718$, $1$; suma $23{,}80$; $p=(0{,}8438;\allowbreak\,0{,}1142;\allowbreak\,0{,}0420)$; $H=0{,}1433+0{,}2478+0{,}1332=0{,}5243$.
  \item $\tau=2$: $z/\tau=(1{,}5;\allowbreak\,0{,}5;\allowbreak\,0)$; $e^{1{,}5}=4{,}482$, $e^{0{,}5}=1{,}649$, $1$; suma $7{,}130$; $p=(0{,}6285;\allowbreak\,0{,}2312;\allowbreak\,0{,}1402)$; $H=0{,}2919+0{,}3386+0{,}2755=0{,}9060$.
  \item Varianza con $\tau=1$: $\mu=3\cdot 0{,}8438+1\cdot 0{,}1142=2{,}646$; $\sum p_iz_i^{2}=9\cdot 0{,}8438+0{,}1142=7{,}708$; $\Var=7{,}708-2{,}646^{2}=7{,}708-6{,}999=0{,}7093$.
  \item Derivada en $b=1$: $dH/db=-1\cdot 0{,}7093=-0{,}7093$.
\end{enumerate}
\resultado \textbf{$H=0{,}107$, $0{,}524$ y $0{,}906$ nats para $\tau=0{,}5$, $1$ y $2$; $\Var_p(z)=0{,}709$ y $dH/db=-0{,}709$ en $b=1$.}
\comprobacion la entropía crece con $\tau$ y queda por debajo de $\ln 3=1{,}099$. La derivada predice, al pasar de $b=1$ a $b=0{,}5$ ($\tau=2$), un aumento de unos $0{,}709\cdot 0{,}5=0{,}35$ nats en primera aproximación; el aumento real es $0{,}38$, mayor porque la varianza crece al disminuir $b$.
\end{solucion}

\begin{solucion}{eje:5-topp}
\datos $q=(0{,}5;\allowbreak\,0{,}2;\allowbreak\,0{,}15;\allowbreak\,0{,}1;\allowbreak\,0{,}05)$, ya ordenada.
\planteamiento \cref{prop:nucleo}: acumular hasta alcanzar $\rho$, conservar esos tokens y renormalizar. La temperatura $\tau=0{,}5$ sobre logits $\ln q_i$ eleva cada $q_i$ al cuadrado antes de normalizar.
\pasos
\begin{enumerate}[label=\arabic*.]
  \item Acumuladas: $0{,}5$; $0{,}7$; $0{,}85$; $0{,}95$; $1$.
  \item $\rho=0{,}8$: la primera acumulada $\ge 0{,}8$ es $0{,}85$, con 3 tokens. Dividiendo por $0{,}85$: $(0{,}5882;\allowbreak\,0{,}2353;\allowbreak\,0{,}1765;\allowbreak\,0;\allowbreak\,0)$.
  \item $\rho=0{,}9$: la primera $\ge 0{,}9$ es $0{,}95$, con 4 tokens. Dividiendo por $0{,}95$: $(0{,}5263;\allowbreak\,0{,}2105;\allowbreak\,0{,}1579;\allowbreak\,0{,}1053;\allowbreak\,0)$.
  \item (b) $q_i^{2}=(0{,}25;\allowbreak\,0{,}04;\allowbreak\,0{,}0225;\allowbreak\,0{,}01;\allowbreak\,0{,}0025)$, de suma $0{,}325$; normalizada: $(0{,}7692;\allowbreak\,0{,}1231;\allowbreak\,0{,}0692;\allowbreak\,0{,}0308;\allowbreak\,0{,}0077)$.
  \item Acumuladas: $0{,}7692$; $0{,}8923\ge 0{,}8$: 2 tokens. Dividiendo por $0{,}8923$: $(0{,}8621;\allowbreak\,0{,}1379;\allowbreak\,0;\allowbreak\,0;\allowbreak\,0)$.
\end{enumerate}
\resultado \textbf{$\rho=0{,}8$: $(0{,}588;\allowbreak\,0{,}235;\allowbreak\,0{,}176;\allowbreak\,0;\allowbreak\,0)$; $\rho=0{,}9$: $(0{,}526;\allowbreak\,0{,}211;\allowbreak\,0{,}158;\allowbreak\,0{,}105;\allowbreak\,0)$; con $\tau=0{,}5$ y $\rho=0{,}8$ quedan dos tokens: $(0{,}862;\allowbreak\,0{,}138;\allowbreak\,0;\allowbreak\,0;\allowbreak\,0)$.}
\comprobacion cada resultado suma~1. Al bajar la temperatura la masa se concentra y el núcleo pasa de 3 a 2 tokens: el número de candidatos se adapta a la forma de la distribución.
\end{solucion}

\begin{solucion}{eje:5-conmuta}
\datos $\tau>0$; logits $z$; para (b), la distribución $q$ del \cref{eje:5-topp} con $\rho=0{,}8$ y $\tau=0{,}5$.
\planteamiento la división por $\tau>0$ conserva el orden de los logits; el núcleo, en cambio, depende de la forma de la distribución.
\pasos
\begin{enumerate}[label=\arabic*.]
  \item (a) Como $\tau>0$, $z_i>z_j$ si y solo si $z_i/\tau>z_j/\tau$, y los empates se conservan; luego $I_k(z/\tau)=I_k(z)$.
  \item Top-$k$ y después temperatura: $\softmax\bigl(\topk(z)/\tau\bigr)$, con $-\infty/\tau=-\infty$. Temperatura y después top-$k$: $\softmax\bigl(\topk(z/\tau)\bigr)$. En ambos casos el resultado es $e^{z_i/\tau}/\sum_{j\in I_k(z)}e^{z_j/\tau}$ en $I_k(z)$ y~0 fuera: coinciden.
  \item (b) Temperatura y después top-$p$: por el \cref{eje:5-topp}, $(0{,}8621;\allowbreak\,0{,}1379;\allowbreak\,0;\allowbreak\,0;\allowbreak\,0)$.
  \item Top-$p$ con $\tau=1$ y después temperatura: el núcleo es $\{1,2,3\}$; dentro de él, $q_i^{2}=(0{,}25;\allowbreak\,0{,}04;\allowbreak\,0{,}0225)$, de suma $0{,}3125$, y la distribución es $(0{,}8;\allowbreak\,0{,}128;\allowbreak\,0{,}072;\allowbreak\,0;\allowbreak\,0)$.
  \item Las dos distribuciones son distintas: incluso el número de tokens con probabilidad positiva difiere (2 frente a 3).
  \item (c) En la \cref{def:topp}, $q$ es la softmax con temperatura: primero se aplica la temperatura y después se calcula el núcleo, como en CS336~\cite[§6]{cs336}.
\end{enumerate}
\resultado \textbf{(a) top-$k$ y temperatura conmutan porque $\tau>0$ no altera el orden de los logits; (b) con top-$p$ no conmutan: $(0{,}862;\allowbreak\,0{,}138;\allowbreak\,0;\allowbreak\,0;\allowbreak\,0)$ frente a $(0{,}8;\allowbreak\,0{,}128;\allowbreak\,0{,}072;\allowbreak\,0;\allowbreak\,0)$; (c) temperatura primero.}
\comprobacion $0{,}8+0{,}128+0{,}072=1$ y $0{,}8621+0{,}1379=1$.
\end{solucion}

\begin{solucion}{eje:5-contexto}
\datos $T=8$ da $0{,}694$ nats y $T=32$ da $0{,}324$; las líneas miden entre unos 58 y 80 caracteres y empiezan por su etiqueta.
\planteamiento identificar qué información útil queda a más de 8 caracteres de la posición que se predice.
\pasos
\begin{enumerate}[label=\arabic*.]
  \item La etiqueta («receta.», «objeto.», «diálogo.», «parte del tiempo.») está al principio de la línea; a partir del noveno carácter queda fuera de la ventana.
  \item Sin la etiqueta, el modelo no sabe qué plantilla escribe: ante «sirve para» solo ve «e\vs para\vs» y no sabe si sigue el uso de un objeto o el final de una receta.
  \item La concordancia de género depende del sustantivo («un reloj negro»), que suele estar a más de 8 caracteres del adjetivo.
  \item Las partes fijas de las plantillas («grados y máxima de») se siguen prediciendo bien: 8 caracteres bastan para reconocerlas.
  \item Exceso de pérdida: $0{,}694-0{,}324=0{,}370$ nats por carácter; en una línea de 58 caracteres, $0{,}370\cdot 58=21{,}5$ nats.
  \item La etiqueta sola aporta $\ln 4=1{,}386$ nats; el exceso es unas 15 veces mayor, porque la incertidumbre sobre la plantilla reaparece en cada elección que depende de ella.
\end{enumerate}
\resultado \textbf{con $T=8$ el modelo pierde la etiqueta y el sustantivo del que dependen las elecciones: las palabras fijas de las plantillas siguen siendo predecibles, pero colores, materiales, usos y concordancias se vuelven inciertos; el exceso, unos $21{,}5$ nats por línea, supera con mucho los $1{,}386$ nats de la etiqueta.}
\comprobacion la pérdida con $T=8$ sigue muy por debajo de la del bigrama ($2{,}035$): un contexto corto ya permite completar palabras.
\end{solucion}
